% Generated draft: requires mathematical and source review.
\documentclass[11pt]{article}
\usepackage[margin=25mm]{geometry}
\usepackage{fontspec}
\setmainfont{Latin Modern Roman}
\usepackage{amsmath,amssymb,mathtools,unicode-math}
\setmathfont{Latin Modern Math}
\usepackage{enumitem,needspace,xurl,hyperref}
\setlength{\parindent}{0pt}
\setlength{\parskip}{5pt}
\setlength{\emergencystretch}{3em}
\newcommand{\N}{\mathbb{N}}
\newcommand{\Z}{\mathbb{Z}}
\newcommand{\Q}{\mathbb{Q}}
\newcommand{\R}{\mathbb{R}}
\newcommand{\C}{\mathbb{C}}
\newcommand{\F}{\mathbb{F}}
\newcommand{\A}{\mathbb{A}}
\newcommand{\PP}{\mathbb{P}}
\newcommand{\E}{\mathbb{E}}
\newcommand{\eps}{\varepsilon}
\newcommand{\dd}{\,\mathrm{d}}
\begin{document}
\section*{UnsolvedMath ID: 800278}
\textbf{Source problem code:} catalogue entry
\subsection{Definitions and mathematical statement}
Write points of \(\C^2\) as \((\lambda,z)\), and let \(\pi(\lambda,z)=\lambda\), \(\mathbb D=\{\lambda\in\C:|\lambda|<1\}\), and \(\mathbb T=\{\lambda\in\C:|\lambda|=1\}\). Interpret the source phrase ``lies over the circle'' as \(Y\subseteq\pi^{-1}(\mathbb T)\). For a compact set \(Y\subseteq\C^2\), define its polynomial hull by
\[
\widehat Y=\left\{w\in\C^2: |p(w)|\leq\sup_{y\in Y}|p(y)|\text{ for every }p\in\C[\lambda,z]\right\},
\]
with the usual convention \(\widehat\varnothing=\varnothing\).

The question is whether the following assertion holds: for every compact \(Y\subseteq\pi^{-1}(\mathbb T)\) and every open set \(U\subseteq\C^2\) satisfying \(\widehat Y\subseteq U\), there exists a Riemann surface
\[
S\subseteq U\cap\pi^{-1}(\mathbb D)
\]
such that \(\pi|_S:S\to\mathbb D\) is a finite-sheeted branched covering of positive degree.

For this formulation, the conventional embedded interpretation is that \(S\) is a one-dimensional complex submanifold of \(\C^2\), and that \(\pi|_S\) is a proper surjective holomorphic map with finite fibers. Locally, at each point of \(S\), the projection has the form \(t\mapsto t^e\) in suitable coordinates, for an integer \(e\geq1\).

These interpretations require review against the source. The supplied text does not specify whether \(S\) must be connected or whether singular analytic curves represented by their normalizations are allowed. It imposes neither nonemptiness of \(Y\) nor \(\pi(Y)=\mathbb T\), and it states no boundary condition relating \(S\) to \(Y\). No such additional hypotheses are included here; whether any were intended remains unresolved.
\subsection{Short English statement}
Does every open neighborhood of the polynomial hull of a compact set lying over the unit circle contain a Riemann surface whose coordinate projection is a finite-sheeted branched covering of the entire unit disc? The supplied statement leaves possible restrictions on the compact set, the surface, and its boundary unspecified.
\subsection{Sources}
\begin{itemize}
\item Local UnsolvedMath registry: permanent ID \(800278\), exact source problem code \texttt{catalogue entry}, title ``Several Complex Variables Problem List --- Problem 78.'' The registry labels the problem open; this is a supplied status, not an independently verified mathematical assessment.
\item Buzzard, Fornaess, and Krantz, \emph{The Several Complex Variables Problem List} (1995), Problem 78. Supplied source URL: \url{https://arxiv.org/abs/math/9509202}. The record reports access on 2026-10-16 and extraction by \texttt{arxiv\_latex}; these are supplied provenance details, not a claim of independent source access.
\item The quoted problem is attributed to H. Alexander and J. Wermer. The source supplies the historical contact \texttt{u22330@uicvm.bitnet} for Alexander and identifies Wermer with Brown University. The registry separately lists Buzzard, Fornaess, and Krantz in its proposer field and gives the year 1995.
\end{itemize}
\end{document}
